\chapter{放行端口并测试连通}
\setchapterheader{第四章\quad 放行端口并测试连通}
\label{chap:firewall}
\begin{intro}
\begin{itemize}
\item 检查并放行 DDS 通信所需的 UDP 端口。
\item 用 ping 和 ARP 表确认到运控板的链路。
\end{itemize}
\end{intro}
在运行控制程序之前，先确认主机防火墙没有拦截 DDS 通信，再从网络层验证到机器人是否可达。

\section{第 8 步：放行 DDS 通信端口}
\label{step:8}
\begin{command}
sudo ufw status verbose
\end{command}

如果显示 \code{Status: active}，执行：

\begin{command}
sudo ufw allow 7400:7610/udp comment 'CycloneDDS RTPS'
sudo ufw allow 22/tcp
sudo ufw reload
sudo ufw status numbered
\end{command}

\begin{expected}
\begin{itemize}
\item 规则列表中出现 \code{7400:7610/udp} 与 \code{22/tcp}。
\item 若防火墙本就未启用，可跳过本步。
\end{itemize}
\end{expected}

\section{第 9 步：测试能否连上机器人}
\label{step:9}
\begin{command}
ping -c 3 192.168.55.102
ip neigh show dev enp0s5
\end{command}

\begin{expected}
\begin{itemize}
\item ping 应有 0\% packet loss。
\item ARP 表里能看到 \code{192.168.55.102} 对应的 MAC 地址。
\end{itemize}
\end{expected}

\begin{faq}{ping 不通}
检查网线、机器人是否上电、双方掩码是否都是 \code{/24}、桥接网卡是否选错。
\end{faq}
